% !TeX program = pdflatex
% Draft section for WILSON_FIRST_RG_STEP_V4.tex: a closed contour on the torus, with corners, which may wind.
% It contains the closed straight line of WILSON_V4_CLOSED_LINE_DRAFT.tex as the case without corners.
% Notation and labels are those of WILSON_FIRST_RG_STEP_V4.tex (SHA-256 prefix 0f859f2d3652).
% Everything between \maketitle and \end{document} is the text of the section; the preamble is that of V4, with
% xr-hyper for its labels (compile V4 first in the same directory, so that its .aux is there).
% State, what was read in V4 for each statement about its proofs, and what is open:
% WILSON_V4_TORUS_CONTOURS_NOTES.md.
\documentclass[11pt]{article}
\usepackage[T1]{fontenc}
\usepackage[utf8]{inputenc}
\usepackage{amsmath,amssymb,amsthm}
\usepackage{mathtools}
\usepackage{mathrsfs}
\usepackage[a4paper,margin=28mm]{geometry}
\usepackage{enumitem}
\usepackage{graphicx}
\usepackage{xr-hyper}
\usepackage[hidelinks]{hyperref}
\externaldocument[V4-][nocite]{../Continuum Limit Manuscripts/WILSON_FIRST_RG_STEP_V4}[../WILSON_FIRST_RG_STEP_V4.pdf]
\newcommand{\dd}{\mathrm{d}}
\newcommand{\Hol}{\operatorname{Hol}}
\newcommand{\tr}{\operatorname{tr}}
\newcommand{\id}{\mathbf 1}
\newcommand{\iso}{\operatorname{iso}}
\newcommand{\sep}{\operatorname{sep}}
\makeatletter
\newcommand*{\bigcdot}{}
\DeclareRobustCommand*{\bigcdot}{%
  \mathbin{\mathpalette\bigcdot@{}}%
}
\newcommand*{\bigcdot@scalefactor}{.7}
\newcommand*{\bigcdot@widthfactor}{1.15}
\newcommand*{\bigcdot@}[2]{%
  \sbox0{$#1\vcenter{}$}%
  \sbox2{$#1\cdot\m@th$}%
  \hbox to \bigcdot@widthfactor\wd2{%
    \hss
    \raise\ht0\hbox{%
      \scalebox{\bigcdot@scalefactor}{%
        \lower\ht0\hbox{$#1\bullet\m@th$}%
      }%
    }%
    \hss
  }%
}
\makeatother
\setlength{\parindent}{0pt}
\setlength{\parskip}{0.6em}
\emergencystretch=3em
\title{Draft: Closed Contours on the Torus}
\author{}
\date{}

\numberwithin{equation}{section}
\newtheorem{theorem}{Theorem}[section]
\newtheorem{proposition}[theorem]{Proposition}
\newtheorem{lemma}[theorem]{Lemma}
\newtheorem{definition}[theorem]{Definition}
\newtheorem{corollary}[theorem]{Corollary}
\theoremstyle{remark}
\newtheorem{remark}[theorem]{Remark}
\theoremstyle{plain}
\renewcommand{\qedsymbol}{\textit{q.e.d.}}
% The section number W belongs to this standalone file only: it keeps the numbers of the draft apart from the
% numbers of V4, which the references to V4 print.
\renewcommand{\thesection}{W}
\begin{document}
\maketitle

\section{Closed Contours on the Torus}\label{sec:torus-contours}

The square-fit condition $K_{\mu}\geqslant 3(n+2)$ of Section~\ref{V4-sec:extent-weights} does two things at once: it fixes the shape of the contour, and it keeps the contour inside a fundamental domain of the torus. In this section we remove both restrictions. The contour is a closed path along coarse bonds which does not meet itself; it may have any number of corners, it need not be planar, and it may wind around the torus. We extract its scalar from the regional reference integral, as Proposition~\ref{V4-prop:regional-square-extraction} does for a square.

What has to be replaced is, first, the elementary geometry of the square. The gas of Section~\ref{V4-sec:marginal-perimeter} and the group sums of Appendix~\ref{V4-app:signed-defects} use it through three counts: the length of the shortest arc of the contour which contains a set of contacts, the number of cubes along an arc, and the number of traversals near a point. For a general contour all three are governed by one number, the largest ratio between the distance of two sites along the contour and their periodic distance, which we call the chord--arc constant; it is $2$ for a square. Since the distance is the periodic one, this constant also measures how close a winding contour comes to its own translates.

The second point is the comparison of the torus with the infinite lattice. A square is compared with its copy on $\mathbb Z^4$, whose connected sum is then expanded into the coefficients of the sides and of the corners. A winding contour has no such copy, and a contour with many corners has no single model either. We therefore compare the density of the torus gas with a model path for each coarse segment separately: the bent path of Section~\ref{V4-sec:marginal-corner} at the corner nearest to the segment, or the straight line if the contour has no corner. A cluster rooted in the segment whose network is shorter than the distance over which the contour coincides with its model, and shorter than a quarter of the shortest period, lifts to a cluster of the model with the same labels, intervals, and closed sources, so that only its covariance marginals differ, while every longer cluster is paid by its extent. The scalar of the contour is then the product of one factor $e^{\kappa_v}$ for each coarse segment and one factor $e^{\zeta_{v,w}}$ for each corner.

\subsection{Contours and Their Unfolding}\label{sec:torus-setting}

We use fine-lattice units. Let $T=\prod_{\mu=1}^4(\mathbb Z/N_{\mu}\mathbb Z)$ have compatible even fine periods $N_{\mu}= LK_{\mu}$ in the sense of Section~\ref{V4-sec:c4-9-2}, and let the standing conditions \eqref{V4-eq:c-standing-conditions} and the collar conditions \eqref{V4-eq:regional-collar-sizes} hold, while the condition $K_{\mu}\geqslant 3(n+2)$ is not imposed. Then $N_{\min}\geqslant 2LMR_1$ by \eqref{V4-eq:c4-9-3} and $LMR_1>4R_0+8$ by \eqref{V4-eq:regional-collar-sizes}, with the range $R_0=4(1+R_*+R_H+L)$ of Section~\ref{V4-sec:regional-identification}, so that
\begin{equation}\label{eq:torus-period}
    N_{\min}>8R_0+16\;.
\end{equation}
We write $\pi_T \colon \mathbb R^4\to\mathbb R^4/\bigoplus_{\mu}N_{\mu}\mathbb Ze_{\mu}$ for the projection, which maps $\mathbb Z^4$ onto $T$, and we call the elements of $\bigoplus_{\mu}N_{\mu}\mathbb Ze_{\mu}$ period vectors, so that a nonzero period vector has maximum norm at least $N_{\min}$. The periodic maximum-norm distance $d_{\infty,T}$ of two points is the least maximum-norm distance of their preimages. Let $\mathscr V=\{\pm e_1,\ldots,\pm e_4\}$ be the set of the eight orientations.

\begin{definition}[Contours]\label{def:torus-contour}
A \emph{contour} is a closed path $C$ of sites $x_0,x_1,\ldots,x_{n_C}= x_0$ of $T$, $n_C\geqslant 3$, with $v_a= x_{a+1}-x_a\in\mathscr V$ for all $a$, the indices being read modulo $n_C$, such that the sites $x_0,\ldots,x_{n_C-1}$ are distinct, the base point $x_0$ is a coarse centre, and every \emph{corner}, that is, every site $x_a$ with $v_a\neq v_{a-1}$, is a coarse centre. Its traversals are the bonds $b_a=(x_a,x_{a+1})$, $a\in\mathbb Z/n_C\mathbb Z$, in this cyclic order; a traversal of orientation $-e_{\mu}$ is a reversed bond and carries the inverse link. We write $\mathscr K(C)$ for the set of the corners, $v_c= v_{a-1}$ and $w_c= v_a$ for the incoming and the outgoing orientation at a corner $c= x_a$, and $d_C$ for the cyclic distance on $\mathbb Z/n_C\mathbb Z$. A \emph{segment} of $C$ is a maximal run of consecutive traversals of one orientation with its end sites; if $C$ has no corner, $C$ itself is its only segment.
\end{definition}

The Wilson function of a contour is $W_{\chi}(C;U)=\chi(\Hol_UC)$, the ordered product being taken along the traversals from the base point. The square $C_n$ is a contour, and so is the closed line through a coarse centre along a coordinate direction, which has no corner.

Two points of a winding contour which are far apart along it may be close on the torus because the contour returns to the neighbourhood of one of its own translates. To keep track of this we work with the preimage. Let $| C|\subset\mathbb R^4/\bigoplus_{\mu}N_{\mu}\mathbb Ze_{\mu}$ be the union of the closed unit intervals of the bonds of $C$. The \emph{unfolding} of $C$ is the set $\widehat C=\pi_T^{-1}(| C|)\subset\mathbb R^4$, which is invariant under the period vectors; its sites are the preimages of the sites of $C$ and its corners the preimages of the corners. A \emph{segment of} $\widehat C$ is a lift of a segment of $C$: a closed interval between two consecutive corners of $\widehat C$ along a lift of the path, or a full line if $C$ has no corner. Two distinct segments of $\widehat C$ are \emph{adjacent} if they have a corner in common.

\begin{definition}[Shape Constants]\label{def:torus-shape}
For a contour $C$, the \emph{shortest segment} $s_{\min}(C)$ is the least number of bonds of a segment of $C$; the \emph{separation} $\sep(C)$ is the least maximum-norm distance between two distinct non-adjacent segments of $\widehat C$; the \emph{corner isolation} and the \emph{chord--arc constant} are
\begin{equation}\label{eq:torus-shape}
    \iso(C)=\min\Bigl(\frac{s_{\min}(C)}{2},\,\sep(C),\,\frac{n_C}{8}\Bigr),
    \qquad
    \Lambda_C=\max\Bigl(2,\,\max_{a\neq a'}\frac{d_C(a,a')}{d_{\infty,T}(x_a,x_{a'})}\Bigr)\;.
\end{equation}
\end{definition}

The letter $\Lambda_C$ is not related to the regions $\Lambda$ or to the spectral bound of \eqref{V4-eq:marginal-covariance-input}. Both distances in \eqref{eq:torus-shape} see the translates of the contour: the separation because the segments of $\widehat C$ include those of all translates, and the chord--arc constant because the distance in its denominator is periodic.

\begin{lemma}[Segments and Corners]\label{lem:torus-segments}
Let $C$ be a contour.
\begin{enumerate}[label=(\alph*)]
    \item $v_a\neq -v_{a-1}$ for every $a$, the orientations $v_c$ and $w_c$ at a corner are perpendicular, and no bond of $T$ is traversed twice, in either direction.
    \item If $C$ has a corner, then it has at least four segments and as many corners as segments, the two end corners of a segment are distinct, and two distinct segments have at most one corner in common.
    \item Every segment has a number of bonds which is a positive multiple of $L$. Hence $n_C= mL$ for an integer $m\geqslant 1$, and $C$ is the concatenation of the straight fine paths of $m$ coarse bonds, which form a closed coarse path; we index the corresponding coarse segments by $j\in\mathbb Z/m\mathbb Z$, the coarse segment $j$ consisting of the traversals $jL,\ldots,(j+1)L-1$, and write $v(j)$ for its orientation.
    \item Two distinct non-adjacent segments of $\widehat C$ are disjoint, and $s_{\min}(C)\geqslant L$, $\sep(C)\geqslant L$ and $\iso(C)\geqslant L/2$.
\end{enumerate}
\end{lemma}

\begin{proof}
(a) If $v_a=-v_{a-1}$, then $x_{a+1}= x_{a-1}$, and since the sites are distinct the indices $a+1$ and $a-1$ agree modulo $n_C$, which is impossible for $n_C\geqslant 3$. Two elements of $\mathscr V$ which are neither equal nor opposite are perpendicular. Two traversals of one bond in the same direction have the same initial site and are therefore the same traversal, and two traversals of one bond in opposite directions would be $b_a$ and $b_{a+1}$ with $x_{a+2}= x_a$, which is excluded in the same way.

(b) The corners and the segments alternate around the contour, so their numbers agree; call the number $k\geqslant 1$. If $k=1$, the orientation changes at one site only, so that all $v_a$ are equal, in contradiction with the definition of a corner. Hence $k\geqslant 2$, and the end corners of a segment, being two sites of $C$ with different indices, are distinct. A segment of $\ell$ bonds along $\pm e_{\mu}$ therefore has $0<\ell< N_{\mu}$, since its $\ell+1$ sites are distinct and differ in the coordinate $\mu$ only. The sum of the $v_a$ is a period vector because the path closes, so for every $\mu$ the signed number of the bonds along $\pm e_{\mu}$ is a multiple of $N_{\mu}$. If exactly one segment were parallel to $e_{\mu}$, its number of bonds would be a multiple of $N_{\mu}$, which we have just excluded; so every axis which carries a segment carries at least two. By (a) at least two axes carry segments, and therefore $k\geqslant 4$. Two distinct segments which have both end corners in common follow each other twice around the contour, which means $k=2$.

(c) If $C$ has corners, the two end corners of a segment are coarse centres, and their lifts along the segment are points of $L\mathbb Z^4$ which differ by $\ell$ times a unit vector, where $\ell$ is the number of bonds; so $L$ divides $\ell$, and the sites of the segment at the distances $L,2L,\ldots$ from its initial corner are coarse centres. The base point is a coarse centre on one of the segments, so its distance from the initial corner of that segment is a multiple of $L$ as well, and the corners are at positions which are multiples of $L$. If $C$ has no corner, it is a closed line through the coarse centre $x_0$ along some $e_{\mu}$, with $n_C= N_{\mu}= LK_{\mu}$ bonds. In both cases the traversals $jL,\ldots,(j+1)L-1$ counted from the base point contain no corner in their interior and join two coarse centres at distance $L$, so they form the straight fine path of a coarse bond.

(d) We first show that two distinct segments $\tilde\sigma_1$, $\tilde\sigma_2$ of $\widehat C$ which meet are adjacent. They are closed axis-parallel intervals, or lines, through sites of $\mathbb Z^4$, with end points in $\mathbb Z^4$; if they meet, they have a site $\tilde x$ in common, because two such sets on a common line which meet share an integer point, and two perpendicular ones meet in a point all of whose coordinates are coordinates of sites. The site $x=\pi_T(\tilde x)$ lies on the segments $\pi_T(\tilde\sigma_1)$ and $\pi_T(\tilde\sigma_2)$ of $C$. If these are distinct, then $x$ is visited once by $C$ and lies on two segments, so it is the corner between them, and $\tilde\sigma_1$, $\tilde\sigma_2$ are the two segments of $\widehat C$ at the corner $\tilde x$, which are adjacent. If they are one segment $\sigma$ of $C$, then $\tilde\sigma_2=\tilde\sigma_1+t$ for a nonzero period vector $t$, and $\tilde x$ and $\tilde x-t$ are two sites of $\tilde\sigma_1$ with the same projection; this is impossible, because the sites of a segment are distinct by (b), or, if $C$ has no corner, because two lines $\tilde\sigma_1\neq\tilde\sigma_1+t$ are parallel and disjoint.

Let now $A$ and $B$ be two disjoint closed axis-parallel intervals of $\mathbb R^4$ with end points in $L\mathbb Z^4$, $A$ parallel to $e_{\mu}$ and $B$ parallel to $e_{\nu}$; we claim that their maximum-norm distance is at least $L$. On each of them all coordinates but one are constant and lie in $L\mathbb Z$. If a coordinate $\lambda\notin\{\mu,\nu\}$ takes different values on $A$ and on $B$, these differ by at least $L$. Otherwise, if $\mu=\nu$, the two intervals lie on one line, where they are disjoint intervals with end points in $L\mathbb Z$; and if $\mu\neq\nu$, let $\beta$ be the value of the coordinate $\mu$ on $B$ and $\alpha$ the value of the coordinate $\nu$ on $A$. The intervals meet if and only if $\beta$ lies in the range of the coordinate $\mu$ on $A$ and $\alpha$ in the range of the coordinate $\nu$ on $B$. One of these fails, say the first, and then $\beta\in L\mathbb Z$ is at distance at least $L$ from that range, whose end points lie in $L\mathbb Z$. This proves the claim. If $C$ has corners, the segments of $\widehat C$ are such intervals, since the corners of $\widehat C$ lie in $L\mathbb Z^4$, the periods being multiples of $L$; two distinct non-adjacent ones are disjoint by the first step, and the claim gives $\sep(C)\geqslant L$. If $C$ has no corner, two distinct lines of $\widehat C$ differ by a nonzero period vector perpendicular to them and are at distance at least $N_{\min}\geqslant L$. Finally $s_{\min}(C)\geqslant L$ by (c), and $n_C\geqslant 4L$ if $C$ has corners by (b) and (c), while $n_C\geqslant N_{\min}$ otherwise.
\end{proof}

\begin{lemma}[Arcs, Cubes, and Counts]\label{lem:torus-counts}
Let $C$ be a contour and $\Lambda=\Lambda_C$.
\begin{enumerate}[label=(\alph*)]
    \item Let $S\subseteq\mathbb Z/n_C\mathbb Z$ be nonempty and let $\delta_S$ be the largest periodic maximum-norm distance of two sites $x_a$, $a\in S$. Then $S$ has at most $2\Lambda\delta_S+1$ elements. If $4\Lambda\delta_S< n_C$, then $S$ lies in a cyclic interval of at most $2\Lambda\delta_S+1$ positions; otherwise $n_C\leqslant 4\Lambda\delta_S$. In both cases the shortest cyclic interval which contains $S$ has at most $4\Lambda\delta_S+1$ positions, and the hull of $S$ in the sense of Section~\ref{V4-sec:joint-extraction} has at most $4\Lambda\delta_S+2$.
    \item Let $\mathcal H$ be a cyclic interval of traversals of $C$ or the whole contour. In the notation of Lemma~\ref{V4-lem:c4-6-env}, $|\pi_1(\mathcal H)|\leqslant 16+16|\mathcal H|/(LM)$.
    \item For every site $x$ of $T$, every integer $s\geqslant 0$ and every $\theta>0$, at most $4\Lambda s+1$ traversals of $C$ have their initial site within the periodic maximum-norm distance $s$ of $x$, and $\sum_{a\in\mathbb Z/n_C\mathbb Z}e^{-\theta d_{\infty,T}(x,x_a)}\leqslant 1+2/(e^{\theta/(2\Lambda)}-1)$.
\end{enumerate}
\end{lemma}

\begin{proof}
(a) Fix $a_0\in S$. By the definition of $\Lambda$ every $a\in S$ has $d_C(a,a_0)\leqslant \Lambda d_{\infty,T}(x_a,x_{a_0})\leqslant \Lambda\delta_S$, so $S$ lies in the set $A$ of the positions at cyclic distance at most $\Lambda\delta_S$ from $a_0$. If $2\Lambda\delta_S< n_C$, this is a cyclic interval of at most $2\Lambda\delta_S+1$ positions; otherwise $S$ has at most $n_C\leqslant 2\Lambda\delta_S$ elements. This gives the first two assertions. In the first case of the dichotomy the shortest interval has at most $2\Lambda\delta_S+1$ positions; the hull is this interval if it has fewer than $n_C/2$ positions, and otherwise it is the whole contour, with $n_C\leqslant 2(2\Lambda\delta_S+1)$. In the second case both have at most $n_C\leqslant 4\Lambda\delta_S$ positions.

(b) We cut the consecutive sites of $\mathcal H$, of which there are at most $|\mathcal H|+1$, into at most $|\mathcal H|/(LM)+1$ runs of at most $LM$ consecutive sites. A lift of such a run is a connected set of at most $LM$ sites of $\mathbb Z^4$, so each of its coordinates takes values in a set of at most $LM$ consecutive integers, which meets at most two of the intervals of $LM$ consecutive integers that define the cubes of $\pi_1$. The run therefore meets at most $2^4$ cubes.

(c) The traversals in question form a set $S$ with $\delta_S\leqslant 2s$, which has at most $4\Lambda s+1$ elements by (a). For the sum, let $a_0$ be a position at which $d_{\infty,T}(x,x_a)$ is least. For every $a$ the triangle inequality gives $d_{\infty,T}(x_{a_0},x_a)\leqslant 2d_{\infty,T}(x,x_a)$, hence $d_{\infty,T}(x,x_a)\geqslant d_C(a_0,a)/(2\Lambda)$, and each value $k\geqslant 1$ of $d_C(a_0,a)$ is taken at most twice. The sum is therefore at most $1+2\sum_{k\geqslant 1}e^{-\theta k/(2\Lambda)}$.
\end{proof}

For the square, $\Lambda_C=2$ and the numbers of (a) are those of Section~\ref{V4-sec:extent-weights}: contacts of diameter $\delta_S$ lie in an arc of at most $4\delta_S+1$ traversals, or else $n_C\leqslant 8\delta_S$. For a general contour they give the bound of a contact interval by the extent with a constant which depends on the chord--arc constant. Indeed, a connected support network of a component $\gamma$ has length at most $R_*| A_{\gamma}|+\sum_{F_{\gamma}}d_T(b,c)\leqslant (1+R_*)\ell(\gamma)$, and each contact is within $R_*$ of it (Section~\ref{V4-sec:extent-weights}), so the contacts have $\delta_S\leqslant 3(1+R_*)\ell(\gamma)$, because $\ell(\gamma)\geqslant 1$. By (a) the interval $h_{\gamma}$ and the hull of the contacts then have at most $4\Lambda\delta_S+2\leqslant 14\Lambda(1+R_*)\ell(\gamma)$ positions. Hence \eqref{V4-eq:marginal-hull-extent} holds for every contour with $\Lambda_C\leqslant \Lambda_{\star}$, with the constant
\begin{equation}\label{eq:torus-hull-constant}
    c_h^{\star}=16\Lambda_{\star}(1+R_*)
\end{equation}
in place of $c_h$; for $\Lambda_{\star}=2$ this is the constant $c_h=32(1+R_*)$ of the square.

The model of a contour near one of its corners is the bent path of Section~\ref{V4-sec:marginal-corner}. Let $C$ have corners. For a corner $\tilde c$ of $\widehat C$ over the corner $c$ of $C$, let $\tilde\sigma_{\mathrm{in}}(\tilde c)$ and $\tilde\sigma_{\mathrm{out}}(\tilde c)$ be the segments of $\widehat C$ which end and begin at $\tilde c$, and let
\begin{equation}\label{eq:torus-bent-path}
    B_{\tilde c}=\{\tilde c-tv_c\, | \,t\in\mathbb Z,\ t\geqslant 0\}\cup\{\tilde c+tw_c\, | \,t\in\mathbb Z,\ t\geqslant 0\}
\end{equation}
be the bent path at $\tilde c$, the translate by $\tilde c\in L\mathbb Z^4$ of the path \eqref{V4-eq:corner-bent-path} for the pair $(v_c,w_c)$; it contains the sites of the two segments. We distribute the coarse segments of $C$ among its corners as Section~\ref{V4-sec:marginal-corner} does for the square: if a segment consists of $n_{\sigma}$ coarse segments, the first $\lceil n_{\sigma}/2\rceil$ of them are \emph{assigned} to its initial corner and the last $\lfloor n_{\sigma}/2\rfloor$ to its final corner. A corner $c$ whose incoming and outgoing segments consist of $n_{\mathrm{in}}$ and $n_{\mathrm{out}}$ coarse segments thus receives the coarse segments which correspond to the indices
\begin{equation}\label{eq:torus-assigned-indices}
    I_c=\bigl\{-\lfloor n_{\mathrm{in}}/2\rfloor,\ldots,\lceil n_{\mathrm{out}}/2\rceil-1\bigr\}
\end{equation}
of the coarse segments of the bent path, and every coarse segment of $C$ is assigned to exactly one corner.

\begin{lemma}[Corner Isolation]\label{lem:torus-corner-isolation}
Let $C$ be a contour with corners, $b$ a traversal of a coarse segment assigned to the corner $c$, $\tilde x$ a lift of its initial site, and $\tilde c$ the corner of $\widehat C$ over $c$ at an end of the segment of $\widehat C$ through $\tilde x$. Every site of $\widehat C$ or of $B_{\tilde c}$ which lies neither on $\tilde\sigma_{\mathrm{in}}(\tilde c)$ nor on $\tilde\sigma_{\mathrm{out}}(\tilde c)$ is at maximum-norm distance at least $\iso(C)-L/2$ from $\tilde x$. If $C$ has no corner, every site of $\widehat C$ which does not lie on the line of $\widehat C$ through a given site $\tilde x$ of $\widehat C$ is at distance at least $\sep(C)\geqslant \iso(C)$ from $\tilde x$.
\end{lemma}

\begin{proof}
The last assertion is the definition of the separation, two distinct lines of $\widehat C$ being non-adjacent. Let $C$ have corners, let $\ell= n_{\mathrm{in}}L$ and $\ell'= n_{\mathrm{out}}L$ be the numbers of bonds of the two segments at $\tilde c$, let $\tilde c'$ be the initial corner of $\tilde\sigma_{\mathrm{in}}(\tilde c)$ and $\tilde c''$ the final corner of $\tilde\sigma_{\mathrm{out}}(\tilde c)$, and write $v= v_c$, $w= w_c$. The segments of $\widehat C$ adjacent to $\tilde\sigma_{\mathrm{in}}(\tilde c)$ are $\tilde\sigma_{\mathrm{out}}(\tilde c)$ and the segment $\tilde\sigma'$ which ends at $\tilde c'$; those adjacent to $\tilde\sigma_{\mathrm{out}}(\tilde c)$ are $\tilde\sigma_{\mathrm{in}}(\tilde c)$ and the segment $\tilde\sigma''$ which begins at $\tilde c''$. Every segment of $\widehat C$ which is neither equal nor adjacent to the one of the two which contains $\tilde x$ is at distance at least $\sep(C)$ from it, so that in each of the two cases below one adjacent segment remains to be considered, besides the sites of the bent path.

Assume first that $b$ lies in one of the last $\lfloor n_{\mathrm{in}}/2\rfloor$ coarse segments of the incoming segment, so that $\tilde x=\tilde c'+\alpha v$ with $\alpha\geqslant \ell-\lfloor n_{\mathrm{in}}/2\rfloor L\geqslant \ell/2\geqslant s_{\min}(C)/2$. The segment $\tilde\sigma'$ is perpendicular to $v$ and contains $\tilde c'$, so all its sites have the coordinate of $\tilde c'$ along $v$, which differs from that of $\tilde x$ by $\alpha$. The sites $\tilde c'-tv$, $t>0$, of the first ray of $B_{\tilde c}$ beyond $\tilde c'$ differ from $\tilde x$ by $\alpha+t$ in the same coordinate, and the sites $\tilde c+tw$, $t>\ell'$, of the second ray beyond $\tilde c''$ differ from $\tilde x$ by $t>\ell'\geqslant s_{\min}(C)$ in the coordinate along $w$, which is constant on $\tilde\sigma_{\mathrm{in}}(\tilde c)$.

Assume now that $b$ lies in one of the first $\lceil n_{\mathrm{out}}/2\rceil$ coarse segments of the outgoing segment, so that $\tilde x=\tilde c+iw$ with $0\leqslant i\leqslant \lceil n_{\mathrm{out}}/2\rceil L-1$. The segment $\tilde\sigma''$ is perpendicular to $w$ and contains $\tilde c''=\tilde c+\ell'w$, so its sites differ from $\tilde x$ by $\ell'-i\geqslant \lfloor n_{\mathrm{out}}/2\rfloor L+1>(\ell'-L)/2\geqslant s_{\min}(C)/2-L/2$ in the coordinate along $w$, and the sites $\tilde c+tw$, $t>\ell'$, of $B_{\tilde c}$ differ from it by more. The sites $\tilde c-tv$, $t>\ell$, of the first ray beyond $\tilde c'$ differ from $\tilde x$ by $t>\ell\geqslant s_{\min}(C)$ in the coordinate along $v$, which is constant on $\tilde\sigma_{\mathrm{out}}(\tilde c)$. In both cases the distances are at least $\min(s_{\min}(C)/2-L/2,\sep(C))\geqslant \iso(C)-L/2$.
\end{proof}

\subsection{The Coefficients and the Scalar of a Contour}\label{sec:torus-coefficients}

Sections~\ref{V4-sec:marginal-perimeter} and~\ref{V4-sec:marginal-corner} define the line coefficient $\kappa_v$ for the orientations $v\in\mathscr V_{\square}$ and the corner logarithm $\zeta_{v,w}$ for perpendicular orientations in $\mathscr V_{\square}$. Their constructions use the straight line through the coarse origin, or the two perpendicular rays \eqref{V4-eq:corner-bent-path}, with the stencils, the primitives and the covariance of the infinite lattice, and they do not use the plane which contains the path. We therefore define, by \eqref{V4-eq:marginal-full-perimeter} and \eqref{V4-eq:corner-local-extraction} with the same proofs and the same constants,
\begin{equation}\label{eq:torus-all-orientations}
    \kappa^{\sigma}_v\quad(v\in\mathscr V),
    \qquad
    \zeta^{\sigma}_{v,w}\quad(v,w\in\mathscr V,\ v\perp w),
\end{equation}
together with the rooted densities $a_{v,w}(i)$ and their straight references $k_{v,w}(i)$ of \eqref{V4-eq:corner-rooted-density}, $i\in\mathbb Z$, so that $|\kappa_v|\leqslant C\tau^2$, $\zeta_{v,w}=\sum_{i\in\mathbb Z}(a_{v,w}(i)-k_{v,w}(i))$ and $| a_{v,w}(i)-k_{v,w}(i)|\leqslant C\tau^2e^{-c| i|}$ by \eqref{V4-eq:corner-localisation-bound}. The \emph{scalar of a contour} $C$ is
\begin{equation}\label{eq:torus-scalar-of-contour}
    \kappa_C=\exp\Bigl(\sum_{j\in\mathbb Z/m\mathbb Z}\kappa^{\sigma}_{v(j)}+\sum_{c\in\mathscr K(C)}\zeta^{\sigma}_{v_c,w_c}\Bigr)\;.
\end{equation}
With $\delta m^{\sigma}_v=-\kappa^{\sigma}_v/(L\varepsilon)$ and $z_{v,w}=e^{\zeta_{v,w}}$ of \eqref{V4-eq:corner-local-extraction} it is the product of $e^{-\delta m^{\sigma}_vP_v(C)}$ over the orientations, $P_v(C)$ being the physical length of the traversals of orientation $v$, and of $z^{\sigma}_{v_c,w_c}$ over the corners. For the square $C_n$ this is the scalar \eqref{V4-eq:extracted-loop-scalar}, by \eqref{V4-eq:marginal-square-mass} and \eqref{V4-eq:corner-square-factor}.

Fix a cutoff prescription $\sigma$. Let $D\subseteq E_{*,T}$ be a set of free bonds which contains every free-input stencil of $C$, in which case we say that $D$ is covered. The regional scalar of the contour is \eqref{V4-eq:regional-square-scalar} with $C$ in place of the square,
\begin{equation}\label{eq:torus-regional-scalar}
    S^{\circ}_{D,C,T}=
    \frac{\displaystyle\int\Theta_D^{\sigma}(B)e^{P_{D,T}(B)}
        \chi\bigl(\Hol_{U_{g_0,T}(i_{D,T}B)}C\bigr)\,\dd\mu_{C_{D,T}}(B)}
    {\displaystyle\int\Theta_D^{\sigma}(B)e^{P_{D,T}(B)}\,\dd\mu_{C_{D,T}}(B)}\;,
\end{equation}
with $C_{D,T}=(i_{D,T}^*H_Ti_{D,T})^{-1}$, the interaction $P_{D,T}$ of \eqref{V4-eq:regional-periodic-weight} and the cutoff $\Theta^{\sigma}_D$ of \eqref{V4-eq:regional-weight}. For a coarse segment $j$ let $K_j$ be the union of the free-input stencils of its fine traversals, and put
\begin{equation}\label{eq:torus-distances}
    d_j=\max\bigl(0,d_T(K_j,D^c)-R_0\bigr),
    \qquad \iota_T(C)=\max\bigl(0,\min\bigl(\iso(C),N_{\min}/4\bigr)-R_0\bigr)\;,
\end{equation}
with complements in $E_{*,T}$, the distance to the empty set being infinite and $e^{-\infty}=0$. The first is the distance of \eqref{V4-eq:regional-square-distances}, and the second reduces to the $\iota_T$ defined there when $\iso(C)\geqslant N_{\min}/4$.

\subsection{The Gas of a Contour}\label{sec:torus-gas}

The constants of the gas depend on the contour through the chord--arc constant. Fix $\Lambda_{\star}\geqslant 2$, and put
\begin{equation}\label{eq:torus-label-constant}
    q_*^{\star}= q_*+4\Lambda_{\star}R_*+1\;,
\end{equation}
with the overlap bound $q_*$ and the range $R_*$ of Section~\ref{V4-sec:marginal-perimeter}. The constants of Section~\ref{V4-sec:marginal-perimeter} and of the proof of Proposition~\ref{V4-prop:regional-square-extraction} are explicit functions of $s_*$, $q_*$, $R_*$, $c_h$ and of the spectral and decay constants of the covariance; we denote by a superscript $\star$ the values which they take when $q_*$ and $c_h$ are replaced by $q_*^{\star}$ and by the constant $c_h^{\star}$ of \eqref{eq:torus-hull-constant}, as Section~\ref{V4-sec:marginal-corner} does for the bent path with $q_*'=2q_*$ and $c_h'=2c_h$. In particular $C_{\mathrm{net}}^{\star}=8(1+R_*)(1+q_*^{\star}+c_h^{\star})$.

\begin{proposition}[The Gas of a Contour]\label{prop:torus-gas}
Assume the setting of Section~\ref{sec:torus-setting}, let $\Lambda_{\star}\geqslant 2$, and let $g_0$ be below the common small-coupling thresholds of the flat and regional expansions of Proposition~\ref{V4-prop:regional-square-extraction} taken with the constants $q_*^{\star}$, $c_h^{\star}$. Let $C$ be a contour with $\Lambda_C\leqslant \Lambda_{\star}$ and let $D$ be covered.
\begin{enumerate}[label=(\roman*)]
    \item The numerator and the denominator of \eqref{eq:torus-regional-scalar} are the sums over the compatible families of vacuum components and webs of the products of their activities, as in \eqref{V4-eq:marginal-exact-gases}, where the webs are formed with the shortest cyclic intervals in $\mathbb Z/n_C\mathbb Z$ and the supports are intersected with $D$.
    \item Let $U(\Gamma)$ be the absolute Ursell coefficient of a cluster $\Gamma$ times the positive majorants of its objects, as in \eqref{V4-eq:regional-square-rooted-majorant}. There are $\eta^{\star}>0$ and $C$, depending on $\Lambda_{\star}$ and not on $C$, $D$, $j$ or the periods, such that for every coarse segment $j$
    \begin{equation}\label{eq:torus-rooted-majorant}
        \sum_{\Gamma\, | \,j\in J_{\Gamma}}U(\Gamma)e^{\eta^{\star}\ell(\Gamma)}\leqslant C\tau^2\;.
    \end{equation}
    \item The scalar \eqref{eq:torus-regional-scalar} is nonzero. Its logarithm $F_{D,C,T}$, continued from zero in the traversal-source polynomial \eqref{V4-eq:regional-wilson-polynomial}, is the absolutely convergent sum of the weights $\Phi_{D,C,T}(\Gamma)$ of the clusters with $J_{\Gamma}\neq\varnothing$, and
    \begin{equation}\label{eq:torus-allocation}
        F_{D,C,T}=\sum_{j\in\mathbb Z/m\mathbb Z}\kappa_{D,C,T}(j),
        \qquad
        \kappa_{D,C,T}(j)=\sum_{\Gamma\, | \,j\in J_{\Gamma}}\frac{\Phi_{D,C,T}(\Gamma)}{| J_{\Gamma}|}\;.
    \end{equation}
\end{enumerate}
\end{proposition}

\begin{proof}
The construction of Section~\ref{V4-sec:marginal-perimeter}, from the closed source differences \eqref{V4-eq:finite-source-closed-differences} to the web activities \eqref{V4-eq:marginal-web-activity}, and its regional form in the proof of Proposition~\ref{V4-prop:regional-square-extraction}, use the contour through four properties, each of which holds for $C$.

First, the contour is closed and its holonomy at the reference field $U_0=\id$ is the identity. This is the case for every closed contour, whether or not it winds around the torus, since $U_0$ is the identity on every bond, and so the differences $G_S$ of \eqref{V4-eq:finite-source-closed-differences} are defined for $C$ and sum to $W_{\chi}(C;U_{g_0})$.

Secondly, no bond is traversed twice (Lemma~\ref{lem:torus-segments}(a)), and at most $q_*^{\star}$ labels meet a bond. Indeed, the primitive labels and the bad-bond labels do not involve the contour and are counted in $q_*$. A traversal whose free-input stencil contains a given bond has its initial site within $R_*$ of the bond, and by Lemma~\ref{lem:torus-counts}(c) there are at most $4\Lambda_{\star}R_*+1$ such traversals. The stencil of a traversal is the projection of the stencil of a bond of $\mathbb Z^4$, on whose $R_*$-neighbourhood $\pi_T$ is injective by \eqref{eq:torus-period}, so the support sizes and the ranges are those of the line.

Thirdly, the interval of a component is bounded by its extent, \eqref{V4-eq:marginal-hull-extent} with the constant $c_h^{\star}$, as was shown after Lemma~\ref{lem:torus-counts}. The conversion \eqref{V4-eq:marginal-hull-root} and the interval tree sums are statements about cyclic intervals of $\mathbb Z/n_C\mathbb Z$, and do not involve the shape of the contour.

Fourthly, when a web interval contains the base position in its interior, the cyclicity of the trace moves the base to a position between two web intervals before the Schur factorisation, as for the square.

The remaining ingredients, namely the counts of the primitive labels, the spectral and decay bounds of the principal marginals of $C_{D,T}$, the density estimate \eqref{V4-eq:marginal-density-majorant}, and the source bounds with the quadratic gain of a single mark, do not involve the contour. With the four properties, the proof of \eqref{V4-eq:marginal-exact-gases} gives (i), and the fifth paragraph of the proof of Proposition~\ref{V4-prop:regional-square-extraction}, in which the reference to the square is the bound \eqref{V4-eq:marginal-hull-extent}, gives \eqref{eq:torus-rooted-majorant} with the constants $\star$. As there, the majorants dominate the cluster weights for the source variables in the polydisc $| t_a|\leqslant 2$, so that the logarithm of the ratio exists on that polydisc, vanishes at zero source, and is the absolutely convergent sum of the cluster weights with $J_{\Gamma}\neq\varnothing$. Its exponential at $t_a=1$ is \eqref{eq:torus-regional-scalar}, which is therefore nonzero. Summing the weights $| J_{\Gamma}|^{-1}$ over the marked segments of a cluster gives one, which is \eqref{eq:torus-allocation}.
\end{proof}

\subsection{Short Clusters Lift to the Model Path}\label{sec:torus-lift}

We now take $D= E_{*,T}$ and write $\Phi_{C,T}$ and $\kappa_{C,T}(j)$ for the weights and the densities of Proposition~\ref{prop:torus-gas} at this carrier, so that $F_{C,T}=\log(N^{\sigma}_{C,T}/D^{\sigma}_T)$ is the logarithm of the flat ratio \eqref{V4-eq:marginal-literal-ratio}. For a cluster $\Gamma$ of this gas, let $\mathscr N(\Gamma)$ be the network of the last paragraph of Section~\ref{V4-sec:extent-weights}: a connected set of bonds of $T$ which contains the carrier of every object of $\Gamma$, the lattice paths which realise its forest edges and attach its marked traversals, and every web interval, and which has at most $C_{\mathrm{net}}^{\star}\ell(\Gamma)$ bonds.

Fix a coarse segment $j$ of $C$ and a lift $\tilde x_j$ of its initial site. Its \emph{model path} $\mathsf M_j$ is the following oriented path of $\mathbb Z^4$ through the lift of the segment at $\tilde x_j$. If $C$ has corners and $j$ is assigned to the corner $c$, it is the bent path $B_{\tilde c}$ of \eqref{eq:torus-bent-path} at the corner $\tilde c$ over $c$ at an end of the segment of $\widehat C$ through $\tilde x_j$, and $i_c(j)\in I_c$ denotes the index of the coarse segment of $B_{\tilde c}$ which is the lift of $j$. If $C$ has no corner, it is the line of $\widehat C$ through $\tilde x_j$. The gas of $\mathsf M_j$ is the gas of Section~\ref{V4-sec:marginal-corner} for the bent path, or of Section~\ref{V4-sec:marginal-perimeter} for the line, translated by a vector of $L\mathbb Z^4$; such translations preserve the block grid, from which the infinite-lattice covariance, the primitives and the stencils are built, and with it the reference construction. We write $\Phi_{\mathsf M_j}$ for its weights, which are $\Phi_{v_c,w_c}$ or $\Phi_{v(j)}$, and we put $C_{\star}=\max(C_{\mathrm{net}}^{\star},2C_{\mathrm{net}})$, which bounds the network constants of the three gases.

\begin{lemma}[Short Clusters Lift to the Model Path]\label{lem:torus-lift}
Assume the setting of Section~\ref{sec:torus-setting} and let $C$ be a contour with $\Lambda_C\leqslant \Lambda_{\star}$. Fix a coarse segment $j$ with its model path $\mathsf M_j$, and let $\tilde\jmath$ be the coarse segment of $\mathsf M_j$ which is the lift of $j$. Call a cluster $\Gamma$ of the gas of $C$ with $j\in J_{\Gamma}$, or a cluster $\widetilde\Gamma$ of the gas of $\mathsf M_j$ with $\tilde\jmath\in J_{\widetilde\Gamma}$, \emph{short} when $C_{\star}\ell+R_0<\min(\iso(C),N_{\min}/4)$ for its extent $\ell$. There is a bijection $\Gamma\mapsto\widetilde\Gamma$ from the short clusters of the first kind onto the short clusters of the second kind with the following properties.
\begin{enumerate}[label=(\alph*)]
    \item $\ell(\widetilde\Gamma)=\ell(\Gamma)$ and $| J_{\widetilde\Gamma}|=| J_{\Gamma}|$.
    \item The objects of $\widetilde\Gamma$ are the lifts of the objects of $\Gamma$, in the same order, with the same labels, forests and incompatibilities. The activity of an object of $\widetilde\Gamma$ is given by the integrals \eqref{V4-eq:marginal-component-functional} and \eqref{V4-eq:marginal-web-activity} of the corresponding object of $\Gamma$ in which every block $C_T(b,c)$ of the covariance marginal is replaced by the block $C_{\infty}(\tilde b,\tilde c)$ at the lifted bonds, and $d(\tilde b,\tilde c)= d_T(b,c)< N_{\min}/4$.
\end{enumerate}
\end{lemma}

\begin{proof}
\emph{The lift.} Let $\Gamma$ be a short cluster with $j\in J_{\Gamma}$, and put $\rho=\min(\iso(C),N_{\min}/4)-R_0$. Its network $\mathscr N(\Gamma)$ is connected and has fewer than $\rho$ bonds, so the argument of Section~\ref{V4-sec:c4-11-4} applies to it: we lift the paths of $\mathscr N(\Gamma)$ from a marked traversal of the coarse segment $j$, which we lift into the lift of $j$ at $\tilde x_j$. Two lifts of one vertex would differ by a nonzero period vector, of maximum norm at least $N_{\min}$, while the two paths which give them have together fewer than $N_{\min}/2$ bonds and therefore end at points whose distance is less than $N_{\min}/2$. The lift $\widetilde{\mathscr N}$ is therefore well defined, $\pi_T$ maps its vertices bijectively onto those of $\mathscr N(\Gamma)$, and any two of its vertices $\tilde x$, $\tilde x'$ are at maximum-norm distance less than $\rho\leqslant N_{\min}/4$. For every nonzero period vector $t$ this gives $\|\tilde x-\tilde x'+t\|_{\infty}>3N_{\min}/4>\|\tilde x-\tilde x'\|_{\infty}$, so that the periodic distance of two vertices of $\mathscr N(\Gamma)$ is the distance of their lifts. In particular $d(\tilde b,\tilde c)= d_T(b,c)$ for any two bonds of the carriers.

\emph{The contour.} A traversal of $C$ in $\mathscr N(\Gamma)$ lifts to a bond of $\widehat C$ both of whose end sites are within the distance $\rho$ of the initial site $\tilde x$ of the lifted root traversal, which lies in the lift of the coarse segment $j$. Since $R_0\geqslant 4L$, we have $\rho<\iso(C)-L/2$, and Lemma~\ref{lem:torus-corner-isolation} shows that both end sites lie on $\tilde\sigma_{\mathrm{in}}(\tilde c)\cup\tilde\sigma_{\mathrm{out}}(\tilde c)$, or on the line of $\widehat C$ through $\tilde x$ if $C$ has no corner. Two sites of this set which are the ends of a bond are consecutive sites of the model path, so every lifted traversal is a traversal of $\mathsf M_j$, with the orientation of the contour. All lifted positions lie in a window of fewer than $\rho+1\leqslant n_C/8$ consecutive positions of $\mathsf M_j$, which projects injectively onto a cyclic interval of fewer than $n_C/2$ positions of $C$, in the order of the contour.

\emph{The objects.} The primitive labels, the bad-bond labels and the forests of $\Gamma$ have their supports in $\mathscr N(\Gamma)$. The primitive functions and their supports agree after lifting, as in the comparison preceding \eqref{V4-eq:corner-finite-shape-difference}, and the lifted marks are traversals of $\mathsf M_j$ with the lifted stencils, the stencil of a traversal being determined by its bond. The components, which are the connected components of the graph of the forest and of the label supports, are therefore the lifts of the components of $\Gamma$. On the torus the interval of a component $\gamma$ is the shortest cyclic interval which contains its marks $S_{\gamma}$. It lies in $\mathscr N(\Gamma)$, so it has fewer than $n_C/2$ positions, and a shortest cyclic interval containing $S_{\gamma}$ with fewer than half of the positions is unique (Section~\ref{V4-sec:joint-extraction}). The image of the ordinary interval of $\mathsf M_j$ from the first to the last lifted mark is a cyclic interval which contains $S_{\gamma}$ and is contained in that shortest interval, because the lift of the latter is an interval of $\mathsf M_j$ which contains the lifted marks. The two therefore coincide, and no tie is resolved by the base. Two such intervals intersect on the torus if and only if their lifts intersect, because both lie in the window. The webs, their intervals, the incompatibilities of the objects and the Ursell coefficient of the cluster are therefore those of the lifted cluster $\widetilde\Gamma$, and $J_{\widetilde\Gamma}$ is the lift of $J_{\Gamma}$, which gives (a) with the equality of the distances.

It remains to compare the closed sources. For a web of $\Gamma$ with the marks $S$, the torus source of \eqref{V4-eq:finite-source-closed-differences} is the alternating sum over $R\subseteq S$ of the character of the product of the links $U_{g_0}(b_a)$, $a\in R$, in the cyclic order from the base, all other traversals carrying the identity, whereas the source of the model path, \eqref{V4-eq:marginal-line-source} on the corresponding arc, is the same sum with the product in the order of the path. If the image of the window does not contain the base in its interior, the two orders agree. Otherwise the first is a cyclic rotation of the second, and the character of a product is invariant under cyclic rotations. The links are the same functions of the lifted fields, since the reconstruction of a link depends on its stencil only. Hence the integrands of \eqref{V4-eq:marginal-component-functional} and \eqref{V4-eq:marginal-web-activity} for an object of $\Gamma$ and for its lift differ only in the covariance marginal, which proves (b).

\emph{The inverse.} Let $\widetilde\Gamma$ be a short cluster of the gas of $\mathsf M_j$ with $\tilde\jmath\in J_{\widetilde\Gamma}$. Its network has at most $2C_{\mathrm{net}}\ell(\widetilde\Gamma)\leqslant C_{\star}\ell(\widetilde\Gamma)<\rho$ bonds, with the $\rho$ of the first step, so $\pi_T$ is injective on it and preserves the distances of its vertices, by the inequality of the first step. Its marks are traversals of $\mathsf M_j$ whose end sites are within the distance $\rho<\iso(C)-L/2$ of a site of the lift of $j$; by Lemma~\ref{lem:torus-corner-isolation}, applied to the sites of $B_{\tilde c}$, these end sites lie on $\tilde\sigma_{\mathrm{in}}(\tilde c)\cup\tilde\sigma_{\mathrm{out}}(\tilde c)$, so the marks are lifts of traversals of $C$; if $C$ has no corner, the model path is a line of $\widehat C$ and this holds for all of its traversals. The projected labels, forests and marks are therefore labels, forests and marks of the gas of $C$, and the projected ordinary interval of a component, which has fewer than $n_C/2$ positions, is the shortest cyclic interval containing its marks: a shorter one would intersect it in a cyclic interval containing the marks (Section~\ref{V4-sec:joint-extraction}), whose lift would be a proper subinterval of the interval from the first to the last mark and would contain both of them, which is impossible. Since $\pi_T$ is injective on the network and on the window, two carriers or two intervals intersect after the projection only if they intersect before it, so that the components, the webs and the incompatibilities are preserved. The projection is therefore a short cluster $\Gamma$ of $C$ with $j\in J_{\Gamma}$ and the same extent, and its lift is $\widetilde\Gamma$, while the projection of the lift of a short cluster is that cluster, so that the two maps are inverse to each other.
\end{proof}

\begin{proposition}[The Densities of a Contour]\label{prop:torus-densities}
Under the hypotheses of Proposition~\ref{prop:torus-gas}, with $g_0$ also below the thresholds of Section~\ref{V4-sec:marginal-corner} for the bent path, there are constants $C,c>0$, depending on $\Lambda_{\star}$ and not on $C$, on $j$ or on the periods, such that for every contour $C$ with $\Lambda_C\leqslant \Lambda_{\star}$ and every coarse segment $j$
\begin{equation}\label{eq:torus-density}
    \bigl|\kappa_{C,T}(j)-a_{v_c,w_c}\bigl(i_c(j)\bigr)\bigr|\leqslant C\tau^2e^{-c\iota_T(C)}\;,
\end{equation}
where $c$ is the corner to which $j$ is assigned; if $C$ has no corner, the same bound holds with $\kappa_{v(j)}$ in place of $a_{v_c,w_c}(i_c(j))$. Consequently
\begin{equation}\label{eq:torus-flat}
    \Bigl| F_{C,T}-\sum_{j\in\mathbb Z/m\mathbb Z}\kappa_{v(j)}-\sum_{c\in\mathscr K(C)}\zeta_{v_c,w_c}\Bigr|
        \leqslant C\tau^2\Bigl[me^{-c\iota_T(C)}+|\mathscr K(C)|e^{-cs_{\min}(C)/(2L)}\Bigr]\;.
\end{equation}
\end{proposition}

\begin{proof}
By the translation invariance of the infinite-lattice construction, $a_{v_c,w_c}(i_c(j))$, or $\kappa_{v(j)}$ if $C$ has no corner, is the sum of $\Phi_{\mathsf M_j}(\widetilde\Gamma)/| J_{\widetilde\Gamma}|$ over the clusters of the gas of $\mathsf M_j$ with $\tilde\jmath\in J_{\widetilde\Gamma}$. We split this sum and the sum \eqref{eq:torus-allocation} for $\kappa_{C,T}(j)$ into their short and their remaining clusters, and we pair the short ones by Lemma~\ref{lem:torus-lift}. The difference of the two densities is then the sum of three terms: the sum of $(\Phi_{C,T}(\Gamma)-\Phi_{\mathsf M_j}(\widetilde\Gamma))/| J_{\Gamma}|$ over the short clusters, the sum of $\Phi_{C,T}(\Gamma)/| J_{\Gamma}|$ over the clusters with $j\in J_{\Gamma}$ which are not short, and minus the sum of $\Phi_{\mathsf M_j}(\widetilde\Gamma)/| J_{\widetilde\Gamma}|$ over the clusters with $\tilde\jmath\in J_{\widetilde\Gamma}$ which are not short.

A cluster which is not short has $\ell\geqslant (\min(\iso(C),N_{\min}/4)-R_0)/C_{\star}$, hence $\ell\geqslant \iota_T(C)/C_{\star}$. The second term is therefore at most $C\tau^2e^{-\eta^{\star}\iota_T(C)/C_{\star}}$ in modulus by \eqref{eq:torus-rooted-majorant}, since the majorants dominate the weights, and the third is at most $C\tau^2e^{-\eta'\iota_T(C)/C_{\star}}$ by the rooted bounds \eqref{V4-eq:corner-bent-tail} and \eqref{V4-eq:marginal-cluster-tail} on the infinite lattice, where $\eta'=\eta/2$ is the exponent of \eqref{V4-eq:corner-bent-tail}.

For a short cluster the two weights differ only through the covariance marginals on its carriers (Lemma~\ref{lem:torus-lift}(b)). The union $A_{\Gamma}$ of the carriers lies in $\mathscr N(\Gamma)$, on which $\pi_T$ is a bijection from the lift, so both $(C_T)_{A_{\Gamma}}$ and $(C_{\infty})_{\widetilde A_{\Gamma}}$ are covariances on the same finite set. Two bonds of it are at lifted distance less than $N_{\min}/4$, and for such bonds the periodisation estimate in the proof of Proposition~\ref{V4-prop:regional-square-extraction} (the display preceding \eqref{V4-eq:regional-square-periodic-comparison}) gives
\begin{equation}\label{eq:torus-covariance}
    \bigl\| C_T(b,c)-C_{\infty}(\tilde b,\tilde c)\bigr\|\leqslant Ce^{-\nu N_{\min}/2}e^{-\nu d(\tilde b,\tilde c)}\;.
\end{equation}
We interpolate linearly between the two covariances on $A_{\Gamma}$. Every interpolated covariance and each of its principal marginals has its spectrum in $[h_L^{-1},\gamma_L^{-1}]$, the common exponential decay, and scalar colour blocks, because both ends have these properties and they are preserved by convex combinations and by principal restriction. The density-direction estimates of the proof of Proposition~\ref{V4-prop:regional-square-extraction} therefore apply with $h_A= e^{-\nu N_{\min}/2}$, as they do there for the square, and give the differentiated cluster bound of the proof of Proposition~\ref{V4-prop:regional-coefficient} in the form
\begin{equation}\label{eq:torus-short-difference}
    \bigl|\Phi_{C,T}(\Gamma)-\Phi_{\mathsf M_j}(\widetilde\Gamma)\bigr|\leqslant C\ell(\Gamma)^2e^{-\nu N_{\min}/2}U(\Gamma)\;.
\end{equation}
Half of the extent weight in \eqref{eq:torus-rooted-majorant} absorbs the factor $\ell(\Gamma)^2$, so that the first term is at most $C\tau^2e^{-\nu N_{\min}/2}$ in modulus, and $N_{\min}/2\geqslant 2\iota_T(C)$. This proves \eqref{eq:torus-density} with $c=\min(\eta'/C_{\star},\eta^{\star}/C_{\star},2\nu)$.

For \eqref{eq:torus-flat} we sum over the coarse segments. If $C$ has no corner, \eqref{eq:torus-allocation} and \eqref{eq:torus-density} give the bound with $|\mathscr K(C)|=0$. Otherwise each coarse segment is assigned to one corner, and $j\mapsto i_c(j)$ maps the coarse segments assigned to $c$ bijectively onto $I_c$, so that by \eqref{eq:torus-allocation} and \eqref{eq:torus-density}
\begin{equation}\label{eq:torus-corner-sum}
    \Bigl| F_{C,T}-\sum_{c\in\mathscr K(C)}\sum_{i\in I_c}a_{v_c,w_c}(i)\Bigr|\leqslant C\tau^2me^{-c\iota_T(C)}\;.
\end{equation}
For a corner $c$ the definition of $\zeta_{v_c,w_c}$ gives
\begin{equation}\label{eq:torus-corner-tail}
    \sum_{i\in I_c}a_{v_c,w_c}(i)=\sum_{i\in I_c}k_{v_c,w_c}(i)+\zeta_{v_c,w_c}-\sum_{i\in\mathbb Z\setminus I_c}\bigl(a_{v_c,w_c}(i)-k_{v_c,w_c}(i)\bigr)\;.
\end{equation}
The first sum on the right is $\lfloor n_{\mathrm{in}}/2\rfloor\kappa_{v_c}+\lceil n_{\mathrm{out}}/2\rceil\kappa_{w_c}$, the sum of $\kappa_{v(j)}$ over the coarse segments assigned to $c$, so that the sum of these terms over the corners is $\sum_j\kappa_{v(j)}$. An index $i\notin I_c$ has $| i|\geqslant \min(\lfloor n_{\mathrm{in}}/2\rfloor+1,\lceil n_{\mathrm{out}}/2\rceil)\geqslant s_{\min}(C)/(2L)$, because both segments consist of at least $s_{\min}(C)/L$ coarse segments, and \eqref{V4-eq:corner-localisation-bound} bounds the last sum of \eqref{eq:torus-corner-tail} by $C\tau^2e^{-cs_{\min}(C)/(2L)}$ after a geometric summation. Adding \eqref{eq:torus-corner-sum} and the sum of \eqref{eq:torus-corner-tail} over the corners gives \eqref{eq:torus-flat}.
\end{proof}

\subsection{The Regional Scalar}\label{sec:torus-regional}

\begin{proposition}[The Regional Scalar of a Contour]\label{prop:torus-scalar}
Assume the setting of Section~\ref{sec:torus-setting}, let $\Lambda_{\star}\geqslant 2$, and let $g_0$ be below the thresholds of Propositions~\ref{prop:torus-gas} and~\ref{prop:torus-densities}. There are constants $C,c>0$, depending on $\Lambda_{\star}$, such that for every contour $C$ with $\Lambda_C\leqslant \Lambda_{\star}$ and every covered $D$ the scalar \eqref{eq:torus-regional-scalar} is nonzero, and with its logarithm continued from zero in the traversal-source polynomial \eqref{V4-eq:regional-wilson-polynomial}
\begin{equation}\label{eq:torus-extraction}
    S^{\circ}_{D,C,T}=\kappa_C\,e^{E^{\circ}_{D,C,T}}\;,
\end{equation}
\begin{equation}\label{eq:torus-error}
    \bigl| E^{\circ}_{D,C,T}\bigr|\leqslant C\tau^2\Bigl[\sum_{j\in\mathbb Z/m\mathbb Z}e^{-cd_j}+me^{-c\iota_T(C)}+|\mathscr K(C)|e^{-cs_{\min}(C)/(2L)}\Bigr]\;,
\end{equation}
with the scalar $\kappa_C$ of \eqref{eq:torus-scalar-of-contour} in the same cutoff prescription. The constants do not depend on $C$, on $D$ or on the periods. As in Proposition~\ref{V4-prop:regional-square-extraction}, when $D$ is the carrier of a covered conditional kernel the collar conditions identify the integral \eqref{eq:torus-regional-scalar} with the conditional reference scalar, and they do not enter the cluster comparison.
\end{proposition}

\begin{proof}
The second and third paragraphs of the proof of Proposition~\ref{V4-prop:regional-square-extraction} identify, for the carrier of a covered conditional kernel, the numerator and the denominator of the regional scalar with those of the conditional reference scalar, and they fix the common domain of the implicit solutions on the integration box. They concern the precision, the primitives and the cutoff on $D$, and they involve the contour through the coverage of its stencils only, so they hold for $C$. Nonvanishing and the logarithm are Proposition~\ref{prop:torus-gas}(iii). We define
\begin{equation}\label{eq:torus-error-definition}
    E^{\circ}_{D,C,T}= F_{D,C,T}-\log\kappa_C
        =\sum_{j\in\mathbb Z/m\mathbb Z}\bigl[\kappa_{D,C,T}(j)-\kappa_{C,T}(j)\bigr]
        +\bigl[F_{C,T}-\log\kappa_C\bigr]\;,
\end{equation}
where $\log\kappa_C$ is the exponent of \eqref{eq:torus-scalar-of-contour} and the second equality is \eqref{eq:torus-allocation} at the carriers $D$ and $E_{*,T}$. The last bracket is bounded by \eqref{eq:torus-flat}.

For the sum we follow the boundary comparison in the proof of Proposition~\ref{V4-prop:regional-square-extraction}, from the periodic resolvent identity to \eqref{V4-eq:regional-square-boundary-comparison}, which is made for each root segment $j$ separately. The bound of $C_{D,T}-(C_T)_{D,D}$ and the two derivative terms of the covariance interpolation do not involve the contour. The contour enters through the network of a cluster rooted in the segment $j$, which has at most $C_{\mathrm{net}}^{\star}\ell(\Gamma)+R_0$ bonds and is anchored within $R_0$ of $K_j$, and through the rooted bound, which is \eqref{eq:torus-rooted-majorant} here. For $d_j>2R_*+2R_0$ we split at $\ell= d_j/(2C_{\mathrm{net}}^{\star})$: a shorter cluster has its original stencils inside $D$, hence the same labels and cyclic intervals in the two expansions, and its carriers at distance at least $d_j/2-R_0$ from $D^c$, so that its covariance derivative is bounded as there; the longer ones are paid by \eqref{eq:torus-rooted-majorant}; and bounded $d_j$ are covered by the same rooted bound after enlarging $C$. This gives $|\kappa_{D,C,T}(j)-\kappa_{C,T}(j)|\leqslant C\tau^2e^{-cd_j}$ for every $j$. Adding the two bounds proves \eqref{eq:torus-error}, and \eqref{eq:torus-extraction} is the exponential of \eqref{eq:torus-error-definition}, every logarithm being the connected source branch.
\end{proof}

\begin{corollary}[The Regional Scalar of the Covered Branch]\label{cor:torus-branch}
Assume the conventions of Section~\ref{V4-sec:c4c-1} with a contour $C$, $\Lambda_C\leqslant \Lambda_{\star}$, in place of $C_n$, with its $m$ coarse segments in place of the $4n$ segments of the square, and without the condition $K_{\mu}\geqslant 3(n+2)$, and the hypotheses of Proposition~\ref{prop:torus-scalar}. On every covered branch $\xi$, with the carrier $I=\Lambda_1^*$, the scalar $S^{\circ}_{\xi,C}$ of \eqref{V4-eq:part1-reference-gas} satisfies
\begin{equation}\label{eq:torus-branch}
    S^{\circ}_{\xi,C}= S^{\circ}_{I,C,T}=\kappa_Ce^{E^{\circ}_{\xi}},
    \qquad \kappa_C\neq0\;,
\end{equation}
where $E^{\circ}_{\xi}= E^{\circ}_{I,C,T}$ is a function of the branch through $I$ alone, and, with the distances \eqref{eq:torus-distances} at $D= I$,
\begin{equation}\label{eq:torus-branch-bound}
    \bigl| E^{\circ}_{\xi}\bigr|\leqslant C\tau^2\Bigl[\sum_{j\in\mathbb Z/m\mathbb Z}e^{-cd_j}+me^{-c\iota_T(C)}+|\mathscr K(C)|e^{-cs_{\min}(C)/(2L)}\Bigr]\leqslant 3C\tau^2m\;.
\end{equation}
\end{corollary}

\begin{proof}
The first equality is Lemma~\ref{V4-lem:c4c-2-1} for $C$. Its proof compares two ratios of integrals over $\mathfrak g^I$ factor by factor, and the factors other than the Wilson factor (the precision, the reconstruction, the primitives and the cutoff) do not involve the contour. For the Wilson factor it uses the reference-gas identity \eqref{V4-eq:part1-reference-gas} and the coverage of every traversal. The proof of \eqref{V4-eq:part1-reference-gas} (Lemma~\ref{V4-lem:schur}) uses the contour through the cyclic list of its traversals, through the intervals of the groups, and through $f^{\circ}_{\xi}(0)=\chi(\id)=1$, which holds for $C$ because the reference background is the identity on every bond. On a covered branch each link of $C$ is its reconstruction from the young variable alone, by the argument of Lemma~\ref{V4-lem:covered}, which treats one contour bond at a time. The second equality and the first bound are Proposition~\ref{prop:torus-scalar} at $D= I$, which is covered on a covered branch and is the carrier $D_+$ of Section~\ref{V4-sec:regional-identification}, and the last bound uses that each exponential is at most one and that $|\mathscr K(C)|\leqslant m$, a segment having at least one coarse segment.
\end{proof}

\subsection{The Contour after the First Step}\label{sec:torus-theorem}

It remains to see what becomes of Theorem~\ref{V4-thm:main-theorem}. Its proof takes the scalar from Proposition~\ref{V4-prop:regional-square-extraction}, which Corollary~\ref{cor:torus-branch} replaces, and for the rest it uses the square only through the counts of Lemma~\ref{lem:torus-counts}, at the places which the proof below lists. The constants which these counts enter are fixed before the coupling; for a contour they depend on a bound $\Lambda_{\star}$ of the chord--arc constant, and so do the thresholds on the coupling.

\begin{theorem}[A Contour after the First Renormalisation Step]\label{thm:torus-contour-theorem}
Let $\Lambda_{\star}\geqslant 2$. Assume (H1), (H2) and (H4) of Section~\ref{V4-sec:main-result}, with the constants and the thresholds of (H2) formed for $\Lambda_{\star}$ as in the proof below, and (H3) without the condition $K_{\mu}\geqslant 3(n+2)$ and without the additional condition for (iii). Let $C$ be a contour with $\Lambda_C\leqslant \Lambda_{\star}$, with $m$ coarse segments and $n_C= mL$ traversals, and let $\kappa_C$ be its scalar \eqref{eq:torus-scalar-of-contour}.
\begin{enumerate}[label=(\roman*)]
    \item (Form.) With the error $E^{\circ}_{\xi}= E^{\circ}_{I,C,T}$ of Corollary~\ref{cor:torus-branch}, the terms $E_{\mathrm{def}}$ of \eqref{V4-eq:c4-1-1} and $E_{\mathrm{bdry}}$ of \eqref{V4-eq:c3-6-2} for $C$, and $\kappa_C\neq0$,
    \begin{equation}\label{eq:torus-theorem-identity}
        \mathbb E_{\varepsilon}[W_{\chi}(C)]=\kappa_C\,\mathbb E_{\mathrm h}\bigl[\mathbf 1_{A_{\mathrm{cov}}}e^{E^{\circ}_{\xi}}W_{\chi}(C;U_1)\bigr]+E_{\mathrm{def}}+E_{\mathrm{bdry}}\;.
    \end{equation}
    \item (Bounds.) On $A_{\mathrm{cov}}$,
    \begin{equation}\label{eq:torus-theorem-bounds}
        \begin{gathered}
            | E^{\circ}_{\xi}|\leqslant C\tau^2\Bigl[\sum_{j\in\mathbb Z/m\mathbb Z}e^{-cd_j}+me^{-c\iota_T(C)}+|\mathscr K(C)|e^{-cs_{\min}(C)/(2L)}\Bigr]\leqslant 3C\tau^2m\;,\\
            | E_{\mathrm{def}}|\leqslant C'_Ln_Ce^{-(c_*/2)\pi_*^2}+e^{C_{\chi} n_C\rho(\varepsilon_1+e^{-cb/(LM)})}-1\;,\\
            | E_{\mathrm{bdry}}|\leqslant C_{\mathrm{resp},\chi}\,n_C(1+g_0r)\,e^{-c_*\pi_*^2}\;.
        \end{gathered}
    \end{equation}
    \item (Size of the leading term.) $|\mathbb E_{\mathrm h}[\mathbf 1_{A_{\mathrm{cov}}}e^{E^{\circ}_{\xi}}W_{\chi}(C;U_1)]|\leqslant e^{3C\tau^2m}$, and
    \begin{equation}\label{eq:torus-theorem-moment}
        \mathbb E_{\mathrm h}\bigl[\mathbf 1_{A_{\mathrm{cov}}}e^{| E^{\circ}_{\xi}|}\bigr]\leqslant \exp\Bigl[C\tau^2\bigl(me^{-c\iota_T(C)}+|\mathscr K(C)|e^{-cs_{\min}(C)/(2L)}\bigr)+n_C\delta^I_q\Bigr]\;,
    \end{equation}
    with $\delta^I_q\leqslant \exp[-c_q\pi_*^2/(2C_K)]$ as in \eqref{V4-eq:c5-7-expbound-I}.
\end{enumerate}
The constants $C$, $c$, $C_{\chi}$ and the rate $\delta^I_q$ depend on $\Lambda_{\star}$ in addition to the quantities named in Section~\ref{V4-sec:main-result}; $C'_L$, $c_*$, $c_q$, $C_K$ and $C_{\mathrm{resp},\chi}$ do not. None of them depends on $C$, on the torus or on the history.
\end{theorem}

\begin{proof}
We go through the proof of Theorem~\ref{V4-thm:main-theorem} and name, at each place at which it uses the square, the statement which takes its place.

\emph{The construction and the identity.} Sections~\ref{V4-sec:sourced-effective-interaction} and~\ref{V4-sec:joint-extraction} use the contour through the function $f_{\xi}(a)= W_{\chi}(C;U_{\xi}(a))$, which is bounded by one and depends on the young variable through the free-input stencils of the contour bonds, and through the cyclic list of the traversals, from which the contact sets, the hulls, the groups and their intervals are formed (Definition~\ref{V4-def:groups}); the Schur identity of Lemma~\ref{V4-lem:schur} uses that the reference background is the identity on every bond, so that $f^{\circ}_{\xi}(0)=1$. All this is available for a contour, which is closed and traverses no bond twice (Lemma~\ref{lem:torus-segments}(a)). The identity \eqref{V4-eq:c4-8-1} is the sum over the branches of \eqref{V4-eq:c-one-step-branch-identity}, with the integrability which Lemma~\ref{V4-lem:c4-2-1} and Section~\ref{V4-sec:c3-6} provide from $| f_{\xi}|\leqslant 1$, and Corollary~\ref{cor:torus-branch} substitutes the scalar; this is (i), as in Corollary~\ref{V4-cor:c4c-5-2}. The first bound of (ii) is \eqref{eq:torus-branch-bound}.

\emph{The exterior response.} Appendix~\ref{V4-app:exterior-chessboard} uses the contour through the set $\Gamma_C$ of its bonds and of their free-input templates, with $|\Gamma_C|\leqslant C_Ln_C$, through the set $I_C$ of its free inputs, and through the multiplicity with which a link is traversed, which is one. Its sums are sums over single bonds of $\Gamma_C$ of lattice sums and of one-point probabilities. Proposition~\ref{V4-prop:c3-7-5} therefore holds for $C$ with its constants, which is the third bound of (ii).

\emph{The signed defects.} Appendix~\ref{V4-app:signed-defects} uses the square at the following places. (1) The geometry paragraph of Section~\ref{V4-sec:c4-6}, the bound of a contact hull by the diameter of the contacts: by Lemma~\ref{lem:torus-counts}(a) the hull has at most $4\Lambda_{\star}\delta_S+2$ traversals. The contacts of a component $\gamma$ lie within $2L$ of its domains, a wall-connected union of at most $n_{\gamma}$ cubes of side $LM$ (Lemma~\ref{V4-lem:c4-6-env}), so $\delta_S\leqslant LMn_{\gamma}+4L\leqslant 2LMn_{\gamma}$ and $| H_{\gamma}|\leqslant (8\Lambda_{\star}+1)LMn_{\gamma}$; this is the first bound of Lemma~\ref{V4-lem:c4-8-2}(b) with $A_{\mathrm h}^{\star}=8\Lambda_{\star}+1$, and it gives $|\mathcal H(W)|\leqslant CLMn(W)$ and the diameter of the comparison box in Section~\ref{V4-sec:c4-8-4}. (2) Lemma~\ref{V4-lem:c4-6-env}(a), the cubes along an interval: by Lemma~\ref{lem:torus-counts}(b), $| X(W)\cup\pi_1(\mathcal H(W))|\leqslant n(W)+16+16|\mathcal H(W)|/(LM)\leqslant 32\mathfrak s(W)$, so that parts (a)--(d) of that lemma hold with $A_E^{\star}= A_F^{\star}=32(2R_{\mathrm e}+1)^4$. (3) The sums over the contour in Lemma~\ref{V4-lem:c4-8-2}(b): the sum of $e^{-\theta_{\mathrm h}d_C(a,a')}$ is a sum over $\mathbb Z/n_C\mathbb Z$ as before; the sum of $e^{-\theta_{\mathrm x}d_T(a',\Delta)}$ over the traversals is at most $e^{\theta_{\mathrm x}LM}(1+2/(e^{\theta_{\mathrm x}/(2\Lambda_{\star})}-1))$ by Lemma~\ref{lem:torus-counts}(c), applied to a site of $\Delta$; and at most $4\Lambda_{\star}(LM+2L)+1$ traversals have stencils which meet a given cube, by the same lemma. (4) The constant $\eta$ is chosen with $1000\eta A_{\mathrm s}^{\star}<\varpi/8$, $A_{\mathrm s}^{\star}=1+A_{\mathrm h}^{\star}$, and the conditions $C\rho LM\leqslant \eta/8$ of Lemmas~\ref{V4-lem:c4-8-3} and~\ref{V4-lem:c4-6-2} are imposed with the constants so obtained. (5) The number of contour links which a selected free label affects (Section~\ref{V4-sec:c4-4-1}) is at most one plus the number of the rows whose stencil contains the label, since no link is traversed twice; this bound does not depend on the contour. With these replacements the proofs of Lemmas~\ref{V4-lem:c4-8-2}, \ref{V4-lem:c4-8-3}, \ref{V4-lem:c4-6-2} and~\ref{V4-lem:c4-6-3} use only cyclic intervals and the sums just bounded. The lifting of the comparison rectangle of a small group (Sections~\ref{V4-sec:c4-8-4} and~\ref{V4-sec:c4-11-4}) uses the bound of the hull by the extent, which is (1), and $N_{\min}>32b$ of \eqref{V4-eq:c4-9-3}, and it does not use the condition $K_{\mu}\geqslant 3(n+2)$; a group whose network winds around the torus belongs to the large-domain class there. Proposition~\ref{V4-prop:c4-12-3} therefore holds for $C$, with $C_{\chi}$ and $c$ depending on $\Lambda_{\star}$, which is the second bound of (ii).

\emph{The leading term.} The first assertion of (iii) follows from the positivity of $\mathbb E_{\mathrm h}$, $\mathbb E_{\mathrm h}[1]=1$, $| W_{\chi}|\leqslant 1$ and the last bound of \eqref{eq:torus-branch-bound}. For \eqref{eq:torus-theorem-moment} we use Lemma~\ref{V4-lem:c5-7-expmoment} at the carrier $I$, that is \eqref{V4-eq:c5-7-expmoment-I} with $\lambda= C\tau^2$. Its proof uses the contour bond by bond, through $\sum_b\varphi_c(b)\leqslant n_CN'_{\mathrm{lat}}(c)$, and at one place through the number of the contour bonds within the distance $s$ of a bond, which is at most $4(2s+1)$ for the square and at most $4\Lambda_{\star}s+1$ for $C$ by Lemma~\ref{lem:torus-counts}(c); this changes the constant $\Phi_c$ which bounds $\varphi_c$, to $(4\Lambda_{\star}R_0'+1)(1+1/c)+4\Lambda_{\star}/c^2$, and nothing else. The joint witness bound of Lemma~\ref{V4-lem:c5-7-joint} concerns marks at single bonds and does not involve the contour. Multiplying the exponential of the two terms of \eqref{eq:torus-theorem-bounds} which do not depend on the history by the exponential moment of $C\tau^2\sum_je^{-cd_j}$ gives \eqref{eq:torus-theorem-moment}, with $\delta^I_q$ formed as in \eqref{V4-eq:c5-7-expbound-I} with the new $\Phi_c$; the bound $\delta^I_q\leqslant \exp[-c_q\pi_*^2/(2C_K)]$ holds below a threshold on $g_0$ which depends on $\Lambda_{\star}$ through $\Phi_c$.

\emph{The hypotheses.} The condition $K_{\mu}\geqslant 3(n+2)$ of (H3) enters the proof of Theorem~\ref{V4-thm:main-theorem} through Proposition~\ref{V4-prop:regional-square-extraction} and through the statements just listed, and the additional condition for (iii) served to bound $4ne^{-c\iota_T}$ by one; the corresponding term is kept in \eqref{eq:torus-theorem-moment}. The thresholds of (H2) are those of the statements used, formed with the constants $q_*^{\star}$, $c_h^{\star}$, $A_{\mathrm h}^{\star}$, $A_E^{\star}$ and $\Phi_c$ above.
\end{proof}

\begin{remark}[The Square and the Closed Line]\label{rem:torus-special-cases}
For the square $C_n$ on a torus with $K_{\mu}\geqslant 3(n+2)$ we have $\Lambda_C=2$, $m=4n$, $s_{\min}(C)= nL$ and $\sep(C)= nL$, the nearest translate being at distance at least $2(n+3)L$, so that $\iso(C)= nL/2$, and \eqref{eq:torus-error} has the form of \eqref{V4-eq:regional-square-error}. For the closed line through a coarse centre along $\pm e_{\mu}$ there is no corner, $m= K_{\mu}$, and the scalar is $e^{K_{\mu}\kappa_v}= e^{-\delta m^{\sigma}_vP(C)}$ with the coefficient of the one orientation $v$ of the line. The square coefficient $\delta m^{\sigma}$ of \eqref{V4-eq:marginal-square-mass} is the mean of the four coefficients $\delta m^{\sigma}_v$, $v\in\mathscr V_{\square}$; Section~\ref{V4-sec:marginal-perimeter} shows that $\kappa_{-v}=\overline{\kappa_v}$ and does not presume any other relation between the coefficients of different orientations, and none is used here.
\end{remark}

\begin{remark}[The Function of the Background]\label{rem:torus-background}
After the scalar is extracted, the function which remains is $W_{\chi}(C;U_1)$. For a contour which winds around the torus it is not close to one when the plaquettes of the background are small, because the contour bounds no surface: the field equal to a fixed $h\in G$ on the bonds of direction $e_1$ and to $\id$ on all other bonds has every plaquette equal to $\id$ and the holonomy $h^{\pm N_1}$ around the closed line along $\pm e_1$, and every element of the compact connected group $G$ is of this form. The bound of Corollary~\ref{V4-cor:main-leading-size}, which sums the plaquettes of the surface bounded by the square, therefore has no counterpart for such a contour.
\end{remark}
\end{document}
